% Options for packages loaded elsewhere
\PassOptionsToPackage{unicode}{hyperref}
\PassOptionsToPackage{hyphens}{url}
\PassOptionsToPackage{dvipsnames,svgnames,x11names}{xcolor}
\documentclass[
  10pt,
  fontset=fandol]{ctexart}
\usepackage{xcolor}
\usepackage[margin=16mm]{geometry}
\usepackage{amsmath,amssymb}
\setcounter{secnumdepth}{-\maxdimen} % remove section numbering
\usepackage{iftex}
\ifPDFTeX
  \usepackage[T1]{fontenc}
  \usepackage[utf8]{inputenc}
  \usepackage{textcomp} % provide euro and other symbols
\else % if luatex or xetex
  \usepackage{unicode-math} % this also loads fontspec
  \defaultfontfeatures{Scale=MatchLowercase}
  \defaultfontfeatures[\rmfamily]{Ligatures=TeX,Scale=1}
\fi
\usepackage{lmodern}
\ifPDFTeX\else
  % xetex/luatex font selection
\fi
% Use upquote if available, for straight quotes in verbatim environments
\IfFileExists{upquote.sty}{\usepackage{upquote}}{}
\IfFileExists{microtype.sty}{% use microtype if available
  \usepackage[]{microtype}
  \UseMicrotypeSet[protrusion]{basicmath} % disable protrusion for tt fonts
}{}
\makeatletter
\@ifundefined{KOMAClassName}{% if non-KOMA class
  \IfFileExists{parskip.sty}{%
    \usepackage{parskip}
  }{% else
    \setlength{\parindent}{0pt}
    \setlength{\parskip}{6pt plus 2pt minus 1pt}}
}{% if KOMA class
  \KOMAoptions{parskip=half}}
\makeatother
\setlength{\emergencystretch}{3em} % prevent overfull lines
\providecommand{\tightlist}{%
  \setlength{\itemsep}{0pt}\setlength{\parskip}{0pt}}
\usepackage{amsmath,amssymb,mathtools,needspace}
\xeCJKDeclareCharClass{CJK}{"2032,"0400->"04FF,"2160->"216B,"2460->"2473,"25A0,"25C7,"25CB,"25CF}
\IfFontExistsTF{Arial Unicode MS}{\xeCJKsetup{AutoFallBack=true}\setCJKfallbackfamilyfont{\CJKrmdefault}{Arial Unicode MS}}{}
\setcounter{secnumdepth}{0}
\setlength{\emergencystretch}{2em}
\allowdisplaybreaks
\usepackage{bookmark}
\IfFileExists{xurl.sty}{\usepackage{xurl}}{} % add URL line breaks if available
\urlstyle{same}
\hypersetup{
  pdftitle={神来之笔},
  colorlinks=true,
  linkcolor={Maroon},
  filecolor={Maroon},
  citecolor={Blue},
  urlcolor={Blue},
  pdfcreator={LaTeX via pandoc}}

\title{神来之笔}
\author{}
\date{}

\begin{document}
\maketitle

\subsection{12.2　神来之笔}\label{ux795eux6765ux4e4bux7b14}

\subsubsection{12.2.1　数学史上一篇千古奇文}\label{ux6570ux5b66ux53f2ux4e0aux4e00ux7bc7ux5343ux53e4ux5947ux6587}

成书于汉代的《九章算术》是我国古代最重要的数学经典。该书``圆田术''给出了

\begin{quote}
校注（agent 补充）：《隋书·律历志》引文中``肭数''\,``盈肭''的``肭''按原页字形保留。
\end{quote}

\href{../../source-images/pdf-296.jpeg}{原页图像}

圆面积的计算公式：

``半周半径相乘得积步。''

即圆面积等于半圆周长与半径的乘积。

这一事实是人们所熟知的，然而由于圆是曲边图形，对古人来说，计算圆面积是个数学难题。

刘徽在注《九章算术》时撰写了《圆田术注》，约 1 800 字，后人称之为《割圆术》。割圆术证明了圆面积公式，并且提供了一个计算圆周率的优秀算法。

\textbf{刘徽的《割圆术》是一篇千古奇文。书中有许多亮点，譬如，早在 1800 年前，刘徽就在人类数学史上首次提出了极限观念。}

刘徽从圆的内接正六边形做起，令边数逐步倍增，计算圆内接正 \(n\) 边形面积 \(S_n\)（\(n=6,12,24,\cdots\)），并建立了圆面积 \(S^*\) 的一个逼近序列

\[
S_6\to S_{12}\to S_{24}\to S_{48}\to\cdots\to S^*.
\]

\textbf{这种加工手续开创了极限计算的先河。}

相比之下，古希腊人畏惧无穷，阿基米德的穷竭法同极限思想毫不相干，因而与刘徽的``割圆术''不可同日而语。

其次，刘徽的割圆术中提出了一种高明的逼近策略，建立了下列\textbf{双侧逼近公式}

\[
S_{2n}<S^*<S_{2n}+(S_{2n}-S_n).
\]

\textbf{在逼近过程中，刘徽直接用数据的偏差 \(S_{2n}-S_n\) 作为校正量，生成圆面积 \(S^*\) 的强近似值，从而舍弃了圆的外切多边形的计算，相比阿基米德的穷竭法显著地节省了计算量。}

最后，特别值得指出的是，\textbf{刘徽的割圆术提出了一种逼近加速技术，在《割圆术》末尾，刘徽突然发力给出了一个神奇的精加工算式------我们称之为``刘徽神算''。}

\textbf{本章推荐的逼近加速技术正是从刘徽神算中感悟并提炼出来的。}

\subsubsection{12.2.2　``一飞冲天''的``刘徽神算''}\label{ux4e00ux98deux51b2ux5929ux7684ux5218ux5fbdux795eux7b97}

前已指出，阿基米德用穷竭法割到内接与外切正 96 边形，获得圆周率 \(\pi=3.14\)，这项成就开创了圆周率科学计算的新纪元。

人们自然会问，阿基米德为什么割到正 96 边形就终止了呢？他为什么不再继续割下去？显然更高精度的圆周率是诱人的。

这个问题的原因很简单，实际计算就会明白，在割圆过程中，每二分一次都要耗费相当大的计算量（对于古人来说，这种计算量是很大的），而且，少数几次割圆对改善精度意义不大。面对这种现实，阿基米德终止于正 96 边形得出圆周率 3.14，这种做法是明智的。

然而刘徽却不满足这个现状。他取圆半径 \(r=10\) 寸（即 1 尺）进行计算，发现正 96 边形二分割圆前后的两个结果

\href{../../source-images/pdf-297.jpeg}{原页图像}

\[
S_{96}=313\frac{584}{625},\quad S_{192}=314\frac{64}{625}
\]

都相当于 \(\pi=3.14\)，它们太粗糙了。面对这种情况，刘徽突发奇想：在几乎不耗费计算量的前提下，能否通过某种简单的加工手续，将两个粗糙的近似值 \(S_{96}\)，\(S_{192}\) 加工成高精度的结果呢？

又想化粗为精，又不愿耗费计算量，这似乎有点异想天开。

《割圆术》下篇一开头，刘徽突然``发力''，他将偏差值

\[
\Delta=S_{192}-S_{96}=\frac{105}{625}
\]

乘以因子

\[
\omega=\frac{36}{105}
\tag{12.2.1}
\]

作为 \(S_{192}\) 的校正量，求得

\[
\begin{aligned}
\widehat S&=S_{192}+\frac{36}{105}(S_{192}-S_{96})\\
&=314\frac{64}{625}+\frac{36}{105}\left(314\frac{64}{625}-313\frac{584}{625}\right)=314\frac{4}{25}.
\end{aligned}
\tag{12.2.2}
\]

\textbf{刘徽指出，这样加工的结果 \(314\dfrac{4}{25}\) 相当于正 3 072 边形的面积。}

\textbf{这样得出的结果 \(S_{3\,072}=314\dfrac{4}{25}\) 相当于圆周率 \(\pi=3.141\,6\)，比 \(\pi=3.14\) 一下子提高了两个数量级。真是一跃千里，一飞冲天。}

我们称这个数学案例为``刘徽神算''。

\textbf{刘徽神算化粗为精的加工效果太神奇了。它的设计机理已大大地超出了人们想象力的限度，因而虽历经千年，至今尚未获得人们的理解和接纳，而一直被禁锢在数学古籍之中。}

该怎样破解刘徽神算的玄机呢？

``探赜索隐，钩深致远。''（《周易·系辞》）本章试图直面这样的课题：探究刘徽神算的隐秘事理，钩取深沉的规律和法则，以破解长期困扰数学界的逼近加速问题。

探索将针对刘徽神算围绕下列三个问题展开：

（1）刘徽神算的校正因子 \(\omega=\dfrac{36}{105}\) 是从哪里来的？

（2）感悟刘徽神算，校正因子应当具有怎样的特质？

（3）推广刘徽神算，归纳出一类普适性的逼近加速法则。

\end{document}
